\chapter{Complexity — Oblivious Ledger}
%%% generated by the blueprint pipeline from TCSlib.Complexity.TuringMachine.Robustness.ObliviousLedger
\section{Overview}

This module closes the quadratic cost accounting for the length-only oblivious
schedule: a purely arithmetic ledger combines the clock-capture, initialization, and
macrostep costs into a uniform quadratic halting bound for the schedule machine
(\lean{Complexity.obliviousSchedule_halts}) and hence for the candidate oblivious
simulator (\lean{Complexity.obliviousCandidate_halts}).  It also provides the
decorated-configuration bridge \lean{Complexity.decoratedCfg} with its initialization
and step laws, used by the data-machine correctness proof.

\section{Declarations}

\begin{lemma}[Arithmetic ledger for the schedule's cost]
\lean{Complexity.obliviousLedger_bound}
\label{Complexity.obliviousLedger_bound}
\leanok
\difficulty{4}
Let $a, b, u, n, \tau, w \in \bbn$ satisfy $n \le u$, $\tau \le b\,(u+1)$, and
$w \le b\,(u+1)$.  Then
\[
  \tau + 3n + 18 + (u+1)\bigl(2w + (a+1) + 4\bigr)
  + 22\,(a+1)(u+1) + 18\,\bigl((a+1)(u+1)\bigr)^2
  \;\le\;
  \bigl(18\,(a+1)^2 + 23\,(a+1) + 3b + 25\bigr)\,(u+1)^2 .
\]
The left-hand side gathers, as abstract summands, the costs of a clock run of length
$\tau$, an input reset over $n$ cells, the fixed overheads, $u+1$ budget rounds over a
word of width $w$, and the allocation and sweep stages measured in the unary budget
$(a+1)(u+1)$; the inequality itself is a statement of pure natural-number arithmetic.
\end{lemma}

\begin{lemma}[Quadratic halting bound for the schedule]
\lean{Complexity.obliviousSchedule_halts}
\label{Complexity.obliviousSchedule_halts}
\leanok
\uses{Complexity.budgetValue, Complexity.budgetValue_bits, Complexity.clockStageCfg, Complexity.clockStageCfg_captures, Complexity.clockStageCfg_setup, Complexity.clockTape, Complexity.oblEmbed, Complexity.obliviousLedger_bound, Complexity.obliviousSchedule, Complexity.setupCfg_finish, Complexity.setupCfg_initializes, Turing.FinTM, Turing.FinTM.ComputesInTime, Turing.FinTM.computesInTime_iff, Turing.MultiTapeTM.initCfg, Turing.MultiTapeTM.output_length_le, Turing.MultiTapeTM.runFrom, Turing.MultiTapeTM.runFrom_add}
\difficulty{5}
Let $W$ be a machine over the Boolean alphabet, let $a, b \in \bbn$, and let
$T \colon \bbn \to \bbn$ satisfy $n \le T(n)$ for every $n$.  Suppose that on every
Boolean input $x$ the machine $W$ halts within $b\,(T(\abs{x})+1)$ steps with output
exactly the list of (little-endian) binary digits of $T(\abs{x})$.  Then for every
Boolean input $x$, the length-only schedule machine built from $W$ and the padding
parameter $a$, started on the embedding of $x$ into the schedule alphabet (each bit
sent to its tagged bit symbol), halts --- through its clock capture, input reset,
copy, fixed-width budget conversion, guide allocation, all macrosteps, and the final
counter test --- within a uniform quadratic budget: there exists a time
\[
  t \;\le\; \bigl(18\,(a+1)^2 + 23\,(a+1) + 3b + 25\bigr)\,\bigl(T(\abs{x})+1\bigr)^2
\]
after which its control state is the halting state.
\end{lemma}

\begin{lemma}[Quadratic halting bound for the candidate]
\lean{Complexity.obliviousCandidate_halts}
\label{Complexity.obliviousCandidate_halts}
\leanok
\uses{Complexity.decorateTM_run, Complexity.oblEmbed, Complexity.obliviousCandidate, Complexity.obliviousDataInit, Complexity.obliviousSchedule, Complexity.obliviousSchedule_halts, Complexity.obliviousSchedule_output, Complexity.obliviousVisit, Complexity.parallelTM_run, Turing.FinTM, Turing.FinTM.ComputesInTime, Turing.MultiTapeTM.initCfg, Turing.MultiTapeTM.runFrom}
\difficulty{4}
Let $W$ and $M$ be machines over the Boolean alphabet, let $a, b \in \bbn$, and let
$T \colon \bbn \to \bbn$ satisfy $n \le T(n)$ for every $n$.  Suppose that on every
Boolean input $x$ the machine $W$ halts within $b\,(T(\abs{x})+1)$ steps with output
exactly the list of (little-endian) binary digits of $T(\abs{x})$.  Then for every
Boolean input $x$, the candidate oblivious simulator built from $W$, $M$, and $a$ ---
the length-only schedule decorated by the data transducer for $M$ and realized over
the binary alphabet by transverse binary blocks --- halts on $x$ within the same
uniform quadratic budget: there exists a time
\[
  t \;\le\; \bigl(18\,(a+1)^2 + 23\,(a+1) + 3b + 25\bigr)\,\bigl(T(\abs{x})+1\bigr)^2
\]
after which its control state is the halting state.  No hypothesis on the computation
of the source machine $M$ is needed, and the bound does not depend on $M$.
\end{lemma}

\begin{definition}[Decorated configuration]
\lean{Complexity.decoratedCfg}
\label{Complexity.decoratedCfg}
\leanok
\uses{Turing.Cfg}
Let $c$ be a configuration of a $k$-tape machine with state set $S$ on an input word
$x$ over an alphabet $A$, let $h \in \{0, \dots, k-1\}$ be a designated head index,
let $d$ be an element of a set $D$ of data states, let $l$ tape contents
$\bbz \to A \cup \{\text{blank}\}$ be given, and let an output word over $A$ be given.
The associated decorated configuration is the configuration of a $(k+l)$-tape machine
with state set $S \times D$ on the same input $x$ whose control state is the state of
$c$ paired with $d$ (halted stays halted), whose input head agrees with $c$, whose
first $k$ work tapes and head positions are those of $c$, whose last $l$ (data) tapes
carry the given contents with every data head at the position of the designated head
$h$ of $c$, and whose output is the given word, maintained separately from the silent
schedule.
\end{definition}

\begin{lemma}[Initial decorated configuration]
\lean{Complexity.decoratedCfg_init}
\label{Complexity.decoratedCfg_init}
\leanok
\uses{Complexity.decorateTM, Complexity.decoratedCfg, Turing.FinTM, Turing.MultiTapeTM.initCfg}
\difficulty{3}
Let $P$ be a machine over an alphabet $A$, decorated by a data transducer with $l$
extra tapes, designated head index $h$, a finite data-state set with initial element
$d_0$, and a visit function.  For every input word $x$, the initial configuration of
the decorated machine on $x$ equals the decorated configuration built from the
initial configuration of $P$ on $x$, the initial data state $d_0$, entirely blank
extra tapes, and the empty output: the data tapes begin blank and aligned with the
schedule's origin.
\end{lemma}

\begin{lemma}[One decorated step]
\lean{Complexity.decoratedCfg_step}
\label{Complexity.decoratedCfg_step}
\leanok
\uses{Complexity.decorateTM, Complexity.decoratedCfg, Complexity.setupWrite, Turing.Action.apply, Turing.Cfg, Turing.Cfg.inputSymbol, Turing.Cfg.workTapeSymbols, Turing.FinTM, Turing.MultiTapeTM.step}
\difficulty{4}
Let $P$ be a machine over an alphabet $A$, decorated by a data transducer with $l$
extra tapes, designated head index $h$, initial data state, and visit function.  Let
$c$ be a configuration of $P$ on an input word $x$ whose control state is a live
state $q$, let $d$ be a data state, and fix contents for the $l$ data tapes and an
output word.  Write $v$ for the value of the visit function on $q$, $d$, the input
symbol scanned by $c$, the $k$ work-tape reads of $c$, and the data-tape symbols at
the position of the designated head $h$; thus $v$ consists of a new data state, one
optional write per data tape, and an optional emitted symbol.  Then one step of the
decorated machine from the decorated configuration equals the decorated configuration
built from one step of $P$ from $c$, the new data state, each data tape updated by
its optional write at the position of the designated head, and the output word
extended by the emitted symbol when one is present.  In particular the decorated
transition consists of exactly the prescribed schedule step together with the local
data visit.
\end{lemma}
